% !TeX program = pdflatex
\documentclass[11pt,a4paper]{article}
\usepackage{../../recursos/latex/estilo-notas}

\begin{document}
\cabecalhoaula{11}{KNN e Árvores de Classificação}{AME §8.3--8.6; ISLP §4.7.6 e cap.~8}

%==============================================================================
\section*{Objetivos da aula}
%==============================================================================
Ao final desta aula o aluno deve ser capaz de:
\begin{itemize}[leftmargin=1.4cm]
  \item adaptar o \textbf{KNN} para classificação (voto da maioria) e estimar
        $\Prob(Y=c\mid\x)$ pelas proporções entre vizinhos;
  \item construir uma \textbf{árvore de classificação} e explicar a predição pela
        \emph{moda};
  \item definir os critérios de impureza \textbf{Gini} e \textbf{entropia} e
        justificar por que se preferem à taxa de erro;
  \item usar \textbf{florestas} e \textbf{boosting} para classificação (voto, $m$,
        AdaBoost);
  \item \textbf{comparar} os classificadores vistos no curso quanto a fronteira,
        interpretabilidade e desempenho.
\end{itemize}

\begin{emsala}
Aula que fecha o Bloco~II reaproveitando dois métodos não paramétricos da Parte~I
(KNN, Aula~04; árvores, Aula~06), agora para classificação. Boa notícia: a estrutura
é a mesma; muda só o ``resumo local'' (média $\to$ moda) e o critério de divisão
(EQM $\to$ impureza). Pergunta de abertura: \emph{``como uma árvore decide a melhor
pergunta a fazer em cada nó quando o alvo é uma classe?''}
\end{emsala}

%==============================================================================
\section{KNN para classificação}
%==============================================================================
Adaptação direta da regressão KNN (Aula~04): para classificar $\x$, olhe seus $k$
vizinhos mais próximos e tome a \textbf{classe mais frequente} (voto da maioria):
\begin{equation}\label{eq:knncls}
  \ghat(\x) = \operatorname{moda}\{Y_i : i\in\mathcal{N}_{\x}\}.
\end{equation}
Equivalentemente, é um classificador \emph{plug-in} (Aula~08): estima-se
\[
  \widehat\Prob(Y=c\mid\x) = \frac{1}{k}\sum_{i\in\mathcal{N}_{\x}}\1{Y_i=c}
\]
(a proporção de vizinhos da classe $c$) e toma-se o $\argmax_c$. Essa estimativa
está, por construção, em $[0,1]$.

\begin{ideia}
$k$ controla o balanço viés--variância, como na regressão: $k=1$ dá fronteira muito
``recortada'' (baixo viés, alta variância); $k$ grande dá fronteira suave (alto
viés, baixa variância). Escolha por validação cruzada (Aula~03).
\end{ideia}

\begin{atencao}
Valem os mesmos cuidados da Aula~04: \emph{padronize} as covariáveis (KNN usa
distância) e lembre da \emph{maldição da dimensionalidade} (Aula~05) --- KNN não
seleciona variáveis e degrada com muitas covariáveis irrelevantes. Em binário, use
$k$ \emph{ímpar} para evitar empates.
\end{atencao}

%==============================================================================
\section{Árvores de classificação}
%==============================================================================
São o análogo das árvores de regressão (Aula~06): partição recursiva do espaço de
covariáveis em regiões $R_1,\dots,R_J$. Duas adaptações:

\subsection{Predição: a moda}
Numa folha $R_k$, prediz-se a \textbf{classe majoritária} (em vez da média):
\begin{equation}\label{eq:arvcls}
  \ghat(\x) = \operatorname{moda}\{Y_i : \X_i\in R_k\},\qquad \x\in R_k,
\end{equation}
e $\widehat\Prob(Y=c\mid\x)=\widehat p_{R_k,c}$, a proporção da classe $c$ na folha.

\subsection{Critério de divisão: impureza}
O EQM não faz sentido com $Y$ qualitativo. Em cada região $R$, seja
$\widehat p_{R,c}$ a proporção da classe $c$. Buscam-se divisões que tornem as
folhas \textbf{puras} (idealmente, uma só classe por folha). Mede-se a impureza por:
\begin{align}
  \text{Índice de Gini:}\quad & G(R) = \sum_{c\in\mathcal{C}} \widehat p_{R,c}\,(1-\widehat p_{R,c}),
  \label{eq:gini}\\
  \text{Entropia:}\quad & D(R) = -\sum_{c\in\mathcal{C}} \widehat p_{R,c}\,\log \widehat p_{R,c}.
  \label{eq:entropia}
\end{align}
Ambas são \emph{mínimas} (zero) quando a folha é pura (algum $\widehat p_{R,c}=1$) e
\emph{máximas} quando as classes estão equilibradas. A cada passo escolhe-se a
divisão que mais reduz a impureza (ponderada pelo tamanho das folhas).

\begin{observacao}[Por que não usar a taxa de erro?]
A taxa de erro de classificação $1-\max_c \widehat p_{R,c}$ também mede impureza,
mas é \textbf{insensível} a mudanças nas proporções que não alterem a classe
majoritária. Gini e entropia são \emph{diferenciáveis} e respondem a qualquer
melhora na pureza --- por isso fazem divisões melhores. (A taxa de erro é boa para a
\emph{poda} final.)
\end{observacao}

\begin{exemplo}[Escolha de divisão por Gini]
Para decidir entre dividir por ``Cardiologia'' ou por ``$>100$ leitos'' (predição:
hospital tem tomógrafo), calcula-se o Gini ponderado de cada partição e escolhe-se a
de \emph{menor} valor (maior pureza). No exemplo do [AME], Gini(Leitos)$\approx
0{,}326 <$ Gini(Cardiologia)$\approx 0{,}472$: divide-se por leitos.
\end{exemplo}

\noindent A construção segue as duas etapas da Aula~06: \emph{cresce-se} uma árvore
grande e depois \emph{poda-se} por custo--complexidade ($\alpha$ por CV). As mesmas
vantagens valem: interpretabilidade, categóricas nativas, seleção automática de
variáveis e captura de interações.

%==============================================================================
\section{\emph{Ensembles} para classificação}
%==============================================================================
Toda a Aula~06 se transfere; as árvores individuais agora classificam:
\begin{description}[leftmargin=2.6cm,style=nextline]
  \item[Bagging / Florestas] treina $B$ árvores em amostras bootstrap e agrega por
        \textbf{voto da maioria} (ou média das probabilidades). Florestas
        descorrelacionam sorteando $m$ variáveis por nó; regra empírica em
        classificação: $m\approx\sqrt{d}$. Robustas a $B$; OOB e importância de
        variáveis ``de graça''.
  \item[Boosting] ajuste sequencial. A variante clássica para classificação é o
        \textbf{AdaBoost}: a cada rodada, reponderam-se as observações dando mais
        peso às mal classificadas, e o classificador final é um voto ponderado dos
        aprendizes fracos. Como em regressão, $B$ grande pode superajustar (use
        \emph{early stopping}); XGBoost é a implementação de referência.
\end{description}

%==============================================================================
\section{Comparando os classificadores do curso}
%==============================================================================
\begin{center}\small
\renewcommand{\arraystretch}{1.3}
\begin{tabular}{@{}lp{3.5cm}p{3.2cm}p{3.6cm}@{}}
\toprule
\textbf{Método} & \textbf{Fronteira} & \textbf{Interpretab.} & \textbf{Observações} \\
\midrule
Logística & linear (no \emph{log-odds}) & alta (coefs.) & dá probabilidades; robusta \\
LDA / QDA & linear / quadrática & média & ótimos sob gaussianidade \\
Bayes ingênuo & depende & média & forte com texto; indep.\ condicional \\
KNN & muito flexível, local & baixa & sensível a escala e dimensão \\
Árvore & ``escada'' (eixos) & \textbf{altíssima} & instável sozinha \\
Florestas / Boosting & muito flexível & baixa & topo de desempenho em dados tabulares \\
SVM (kernel) & linear ou não linear & baixa & ótima com classes separadas \\
\bottomrule
\end{tabular}
\end{center}

\begin{ideia}
Não há método universalmente melhor (Aula~05: ``não existe almoço grátis''). A
escolha depende da \emph{estrutura} dos dados e do \emph{objetivo}: precisa de
probabilidades calibradas? interpretabilidade? muitas covariáveis irrelevantes?
fronteira não linear? Compare alguns candidatos por validação cruzada e use as
métricas adequadas (Aula~09).
\end{ideia}

\begin{emsala}
Use o notebook \texttt{Comparação entre classificadores paramétricos} e o
\texttt{Exemplo - KNN e árvores em dados artificiais} para \emph{visualizar} as
fronteiras de decisão lado a lado. É o melhor jeito de fixar a intuição de
``flexível vs.\ rígido'' que percorreu o curso inteiro.
\end{emsala}

%==============================================================================
\section{Prática em Python}
%==============================================================================
\begin{lstlisting}
from sklearn.neighbors import KNeighborsClassifier
from sklearn.tree import DecisionTreeClassifier
from sklearn.ensemble import RandomForestClassifier, AdaBoostClassifier
from sklearn.preprocessing import StandardScaler
from sklearn.pipeline import make_pipeline

knn  = make_pipeline(StandardScaler(),
                     KNeighborsClassifier(n_neighbors=11))   # k por CV
arv  = DecisionTreeClassifier(criterion="gini", ccp_alpha=0.0)  # ou "entropy"
rf   = RandomForestClassifier(n_estimators=500, max_features="sqrt",
                              oob_score=True, random_state=0)
ada  = AdaBoostClassifier(n_estimators=200, learning_rate=0.5)

for nome, modelo in [("KNN", knn), ("arvore", arv), ("RF", rf), ("Ada", ada)]:
    modelo.fit(X_tr, y_tr)
    print(nome, modelo.score(X_te, y_te))   # use tambem as metricas da Aula 09
\end{lstlisting}

%==============================================================================
\section*{Resumo}
%==============================================================================
\begin{itemize}[leftmargin=1.4cm]
  \item \textbf{KNN classificação} \eqref{eq:knncls}: voto da maioria; estima
        $\Prob(Y=c\mid\x)$ por proporções; $k$ por CV; padronize.
  \item \textbf{Árvore de classificação}: prediz a moda \eqref{eq:arvcls}; divide
        por \textbf{Gini} \eqref{eq:gini} ou \textbf{entropia} \eqref{eq:entropia}
        (melhores que a taxa de erro); cresce e poda como na Aula~06.
  \item \textbf{Florestas} ($m\approx\sqrt d$) e \textbf{boosting/AdaBoost} agregam
        árvores por voto --- topo de desempenho em dados tabulares.
  \item Não há classificador universalmente melhor: escolha pela estrutura/objetivo
        e compare por CV com as métricas da Aula~09.
\end{itemize}

\paragraph{Leitura recomendada.} [AME] §8.3 (árvores de classificação, índice de
Gini), §8.4--8.5 (\emph{bagging}, florestas, \emph{boosting}), §8.6 (KNN); [ISLP]
§4.7.6 (KNN) e Capítulo~8 (árvores e \emph{ensembles}, incluindo critérios de
impureza e \emph{boosting}).

\end{document}
